% !TeX root = ../navier-stokes-manuscript.tex
% ============================================================
% 2.5 Simultaneous Sobolev Intersections
% ============================================================

\subsection{Simultaneous Sobolev Intersections}\label{sub:SimultaneousSobolevIntersections}

The spatial column discovered inside critical height has no final finite rung.
Yet any particular derivative asks for only a finite amount of that column.
Suppose, for example, that one wants \(\partial_x^\alpha u\) to be continuous for every \(\lvert\alpha\rvert\le k\).
In three dimensions, the Sobolev embedding
\[
  H^m(\R^3)
  \hookrightarrow
  C^k(\R^3)
\]
holds whenever
\[
  m>k+\frac32.
\]
One requested spatial order \(k\) comes from one finite height \(m\) above its embedding threshold.
As \(k\) ranges through all finite orders, no last \(m\) can be chosen, but no infinite-order norm is required either.
The exact spatial object is the coordinatewise intersection
\[
  \bigcap_{m=0}^{\infty}H^m(\R^3),
\]
whose membership means that every finite Sobolev norm is finite on the same field.
For each fixed \(k\), one of those finite memberships yields the desired \(C^k\) control, so
\[
  \bigcap_{m=0}^{\infty}H^m(\R^3)
  \hookrightarrow
  C^\infty(\R^3).
\]

Finite-time response asks for the same construction in the temporal direction.
Let \(I=(0,T)\) with \(T<T_*\), and let \(X\) be a Hilbert space.
The vector-valued embedding
\[
  H^r(I;X)
  \hookrightarrow
  C^q(\overline I;X)
\]
holds whenever
\[
  r>q+\frac12.
\]
One requested time derivative of order \(q\) comes from one finite temporal height \(r\).
Allowing every finite request gives
\[
  \bigcap_{r=0}^{\infty}H^r(I;X)
  \hookrightarrow
  C^\infty(\overline I;X).
\]
The time intersection is again coordinatewise: each \(q\) chooses its own finite \(r\).

The mixed jet requires these choices at the same time.
Ask for one concrete coordinate,
\[
  \partial_t^2\partial_x^\alpha u,
  \qquad
  \lvert\alpha\rvert\le4.
\]
The single finite membership
\[
  u\in H^3\bigl(I;H^6(\R^3)\bigr)
\]
places that coordinate in
\[
  C^2\bigl(\overline I;C^4(\R^3)\bigr),
\]
because
\[
  3>2+\frac12,
  \qquad
  6>4+\frac32.
\]
Nothing special happened at the pair \((2,4)\).
For arbitrary finite temporal and spatial orders \(q\) and \(k\), choose finite indices satisfying
\[
  r>q+\frac12,
  \qquad
  m>k+\frac32.
\]
Then
\[
  H^r\bigl(I;H^m(\R^3)\bigr)
  \hookrightarrow
  C^q\bigl(\overline I;C^k(\R^3)\bigr).
\]

This repeated finite choice determines the whole velocity target:
\[
  \mathscr V_\infty(I)
  :=
  \bigcap_{r=0}^{\infty}
  \bigcap_{m=0}^{\infty}
  H^r\bigl(I;H^m(\R^3)\bigr).
\]
Its meaning is
\[
  u\in\mathscr V_\infty(I)
  \quad\Longleftrightarrow\quad
  \norm{u}_{H^r(I;H^m)}
  <
  \infty
  \quad
  \text{for every finite pair }(r,m).
\]
The constant may depend on \(r\) and \(m\).
Intersection membership asserts all finite coordinates on the same \(u\); it does not assert one constant uniform across derivative orders.
The embeddings above then give
\[
  \mathscr V_\infty(I)
  \hookrightarrow
  C^\infty\bigl(\overline I\times\R^3\bigr).
\]

The complete intersection also fixes the logical place of critical height.
Choose a temporal index above \(1/2\) and a spatial index above \(3/2\).
The resulting continuous Sobolev trace contains both
\[
  u(t)\in H^{1/2}(\R^3)
  \qquad\text{and}\qquad
  u(t)\in H^{3/2}(\R^3)
\]
at each time in the closed interval.
Only after those rows are supplied do the critical readouts
\[
  H(t)=\frac12\norm{\Lambda^{1/2}u(t)}_2^2,
  \qquad
  E_{3/2}(t)=\frac12\norm{\Lambda^{3/2}u(t)}_2^2
\]
become available.
Critical height motivated the search for the column, while the all-row velocity intersection supplies the critical coordinates used at the endpoint.

The velocity intersection does not contain pressure as one of its elements.
Pressure remains attached through the original recurrence,
\[
  Q_n=\partial_t^np,
\]
the differentiated momentum equation, incompressibility, and the pressure Poisson equation at every finite depth.
Those relations keep the velocity intersection pressure-complete without confusing the pressure jet reconstructed from (u) with membership in the velocity space itself.

The intersection ranges over every finite temporal--spatial pair, while any one estimate controls a finite block.
The adjacent energy levels around a common time trace determine the shape of that block.
